\section{本章小结}
\label{sec:ivp-summary}

本章从规律和起点出发，确定了初值问题的提法。一个初值问题由右端$\bm{f}$、初始时刻$t_0$和初值$\bm{x}_0$组成，答案是一条曲线及其存在区间，两者都要回答。相同的规律从不同的初值出发是不同的曲线，这是方向场所显示的基本事实。初值集中在一个时刻，因此可以沿时间往前推；两点边值的条件分在两头，同样的方程可能无解或有无穷多解。

状态是此刻足以决定下一瞬间变化的量。高阶方程把直到$m-1$阶的导数放进状态，降为一阶系统；单摆需要角度和角速度。时间、状态、外部输入和参数各有角色。右端不显含时间的方程称为自治方程，它的解可以平移；非自治方程可以通过把时间并入状态写成自治方程，代价是增加一维，并把时间上的不正则转移到状态分量上。

解的概念决定了什么才算求解。经典解要在一个区间上处处可微、满足方程和初值，这个要求排除了隐式公式的错误分支，也说明了存在区间可能因定义域的边界或有限时间爆炸而有限。平衡解容易在除法中丢掉，需要单独列出。

解析方法只对具体形式的方程有效。可分离方程由$B(x)=A(t)$得到解，$B$的值域决定存在区间，$b$的零点需要分别处理，$1/b$的可积性决定解能否到达零点或无穷远。齐次型方程和Bernoulli方程用代换化为可分离方程或线性方程，每次代换都要核对可逆范围、零点和取值范围，形式上干净的结果不一定是原问题的全部解。

积分形式把方程写成初值加累积，在连续右端下与微分形式等价。它自带初值，解是自洽的曲线，积分对小扰动不敏感。这个形式让我们能够处理时间方向不连续的可测输入：解的概念放宽为局部绝对连续的曲线，方程几乎处处成立，等价于积分方程处处成立。下表对比了两种解。

\begin{BookInlineTable}
 \centering
 \caption{经典解与绝对连续解}
 \label{tab:ivp-solution-comparison}
 \begin{tabular}{lll}
 \toprule
 & 经典解 & 绝对连续解 \\
 \midrule
 右端 & 连续 & 对时间可测 \\
 曲线 & 连续可微 & 局部绝对连续 \\
 方程 & 处处成立 & 几乎处处成立 \\
 积分方程 & 处处成立 & 处处成立 \\
 例子 & Logistic解 & 切换输入的解 \\
 \bottomrule
 \end{tabular}
\end{BookInlineTable}

没有公式时，仍然要回答解是否存在、是否唯一、能延伸多远，这是第~\ref{chap:ode-existence-continuation}章的任务；解对初值和参数有多敏感，由第~\ref{chap:ode-sensitivity}章估计；怎样算出解的数值，由第~\ref{chap:ode-numerical-methods}章讨论。线性方程可以写出完整的传播公式，是第~\ref{chap:linear-differential-systems}章的内容。
